% Clase 15 --- Entrada oblicua: la hélice (§1.4).
% No hay física nueva acá, sólo una descomposición: la componente paralela a B
% no siente fuerza y sigue de largo, la perpendicular hace el círculo de
% siempre. Lo que la figura tiene que mostrar es que las dos cosas ocurren a la
% vez y de manera independiente.
\begin{center}
\begin{tikzpicture}[scale=1]

  % ---- eje del campo ----
  \draw[figvecaux, figgray, line width=1pt,
        -{Latex[length=2.4mm,width=1.7mm]}] (-3.4,0) -- (3.6,0);
  \node[figlbl, anchor=west] at (3.7,0) {$\vec B$};

  % ---- la hélice ----
  \draw[figred, line width=1pt] plot[domain=0:1080, samples=200, smooth]
    ({-3.0 + \x/180}, {0.85*sin(\x)});

  \node[figcarga, fill=figred, minimum size=6pt] at (-3.0,0) {};

  % ---- descomposición de la velocidad ----
  \draw[figvecres] (-3.0,0) -- (-1.85,0.95);
  \node[figlbl, figred, anchor=south east] at (-1.9,0.95) {$\vec v$};
  \draw[figvecaux, figblue, line width=0.9pt,
        -{Latex[length=2.2mm,width=1.6mm]}] (-3.0,0) -- (-1.85,0);
  \node[figlbl, figblue, anchor=north] at (-2.4,-0.12) {$\vec v_\parallel$};
  \draw[figvecaux, figblue, line width=0.9pt,
        -{Latex[length=2.2mm,width=1.6mm]}] (-3.0,0) -- (-3.0,0.95);
  \node[figlbl, figblue, anchor=east] at (-3.06,0.5) {$\vec v_\perp$};

  % ---- rótulos de las dos partes ----
  \node[figlbl, figblue, anchor=south, align=center] at (0.3,1.6)
    {$\vec v_\perp$: círculo de radio $r = \dfrac{m\,v_\perp}{|q|\,B}$};
  \node[figlbl, figblue, anchor=north, align=center] at (0.3,-1.5)
    {$\vec v_\parallel$: no siente fuerza, sigue de largo};

\end{tikzpicture}\\[0.5em]
{\small\itshape La descomposición funciona porque las dos partes no se estorban.
La componente paralela a $\vec B$ no aparece en el producto vectorial, así que
sobre ella no hay fuerza y se conserva: es un movimiento rectilíneo uniforme a
lo largo del campo. La componente perpendicular hace exactamente el círculo de
la sección anterior, con la salvedad de que en el radio entra $v_\perp$ y no la
rapidez total. Superponer las dos da la hélice ---el paso del tirabuzón es lo
que avanza $\vec v_\parallel$ en un período $T$---. Notar que el período sigue
sin depender de nada de esto: es el mismo $2\pi m/(|q|B)$, así que una entrada
más oblicua no cambia el ritmo de giro, sólo estira la hélice.}
\end{center}
